% =====================================================================
%  minggu-13.tex  --  Slide Pertemuan Minggu ke-13
%  Pengantar Kriptografi (MPFA014) -- Prodi Matematika FMIPA UBB
%
%  Topik silabus : Kriptografi kurva eliptik
%  Sumber materi : buku_ajar/chapters/14 (Bab 14, Bagian 1--4)
%  Sub-CPMK      : 1.4
%
%  Kompilasi dari dalam folder ppt/:
%      pdflatex -interaction=nonstopmode minggu-13.tex   (dua kali)
% =====================================================================

\documentclass[aspectratio=169,11pt]{beamer}
% =====================================================================
%  beamer-preamble.tex  --  prasarana dan gaya slide kuliah
%
%  Mata kuliah : Pengantar Kriptografi (MPFA014)
%  Program     : Program Studi Matematika, FMIPA, Universitas Bale Bandung
%
%  Berkas ini TIDAK memuat \documentclass. Setiap berkas minggu memuatnya
%  sendiri lalu menarik berkas ini:
%
%      \documentclass[aspectratio=169,11pt]{beamer}
%      % =====================================================================
%  beamer-preamble.tex  --  prasarana dan gaya slide kuliah
%
%  Mata kuliah : Pengantar Kriptografi (MPFA014)
%  Program     : Program Studi Matematika, FMIPA, Universitas Bale Bandung
%
%  Berkas ini TIDAK memuat \documentclass. Setiap berkas minggu memuatnya
%  sendiri lalu menarik berkas ini:
%
%      \documentclass[aspectratio=169,11pt]{beamer}
%      % =====================================================================
%  beamer-preamble.tex  --  prasarana dan gaya slide kuliah
%
%  Mata kuliah : Pengantar Kriptografi (MPFA014)
%  Program     : Program Studi Matematika, FMIPA, Universitas Bale Bandung
%
%  Berkas ini TIDAK memuat \documentclass. Setiap berkas minggu memuatnya
%  sendiri lalu menarik berkas ini:
%
%      \documentclass[aspectratio=169,11pt]{beamer}
%      \input{beamer-preamble}
%      \title{...} \subtitle{...} \setWeekInfo{Minggu n}{Topik}
%
%  Berkas minggu hanya berisi \setWeekInfo dan lingkungan frame.
%  Kompilasi dijalankan DARI DALAM folder ppt/ (lihat compile-ppt.bat),
%  karena gambar buku dirujuk dengan jalur relatif ../buku_ajar/.
% =====================================================================

% ---------------------------------------------------------------- bahasa
\usepackage[utf8]{inputenc}
\usepackage[T1]{fontenc}
\usepackage[indonesian]{babel}
\usepackage{lmodern}
\usepackage{microtype}

% ------------------------------------------------------------- matematika
\usepackage{amsmath,amssymb,mathtools}

% ----------------------------------------------------------------- tabel
\usepackage{booktabs}
\usepackage{array}
\usepackage{tabularx}

% --------------------------------------------------------------- grafik
\usepackage{tikz}
\usetikzlibrary{arrows.meta,positioning,shapes.geometric,calc,fit,backgrounds}

% ---------------------------------------------------------- kode program
\usepackage{listings}
\lstdefinestyle{slidecode}{
    basicstyle=\ttfamily\scriptsize,
    keywordstyle=\color{ubbblue}\bfseries,
    commentstyle=\color{ubbgreen}\itshape,
    stringstyle=\color{ubbred},
    numbers=left,
    numberstyle=\tiny\color{gray},
    numbersep=6pt,
    frame=single,
    rulecolor=\color{ubbblue!40},
    backgroundcolor=\color{ubbblue!3},
    breaklines=true,
    showstringspaces=false,
    tabsize=4,
    columns=fullflexible,
    keepspaces=true,
}
\lstset{style=slidecode}

% =====================================================================
%  WARNA INSTITUSI
% =====================================================================
\definecolor{ubbblue}{RGB}{12, 62, 120}      % warna utama
\definecolor{ubbteal}{RGB}{0, 118, 110}      % aksen kedua
\definecolor{ubbgreen}{RGB}{22, 122, 74}     % untuk blok contoh
\definecolor{ubbred}{RGB}{168, 42, 42}       % untuk blok peringatan
\definecolor{ubbgray}{RGB}{70, 76, 84}       % teks sekunder

% =====================================================================
%  IDENTITAS MATA KULIAH
% =====================================================================
\newcommand{\CourseName}{Pengantar Kriptografi}
\newcommand{\CourseCode}{MPFA014}
\newcommand{\CourseSemester}{Semester Ganjil}
\newcommand{\CourseProgram}{Program Studi Matematika}
\newcommand{\CourseFaculty}{Fakultas Matematika dan Ilmu Pengetahuan Alam}
\newcommand{\CourseInstitution}{Universitas Bale Bandung}
\newcommand{\CourseLecturer}{Cecep Suwanda, S.Si., M.Kom.}

% Label minggu, disetel tiap berkas lewat \setWeekInfo
\newcommand{\WeekLabel}{}
\newcommand{\WeekTopic}{}
\newcommand{\setWeekInfo}[2]{%
    \renewcommand{\WeekLabel}{#1}%
    \renewcommand{\WeekTopic}{#2}%
}

% =====================================================================
%  TEMA: header, footer, warna blok
% =====================================================================
\usetheme{default}
\setbeamertemplate{navigation symbols}{}
\setbeamertemplate{blocks}[rounded][shadow=false]

\setbeamercolor{palette primary}{bg=ubbblue, fg=white}
\setbeamercolor{palette secondary}{bg=ubbgray, fg=white}
\setbeamercolor{palette tertiary}{bg=ubbteal, fg=white}

\setbeamercolor{structure}{fg=ubbblue}
\setbeamercolor{normal text}{fg=black!92}
\setbeamercolor{alerted text}{fg=ubbred}

\setbeamercolor{block title}{bg=ubbblue, fg=white}
\setbeamercolor{block body}{bg=ubbblue!6, fg=black!92}
\setbeamercolor{block title alerted}{bg=ubbred, fg=white}
\setbeamercolor{block body alerted}{bg=ubbred!6, fg=black!92}
\setbeamercolor{block title example}{bg=ubbgreen, fg=white}
\setbeamercolor{block body example}{bg=ubbgreen!6, fg=black!92}

\setbeamercolor{itemize item}{fg=ubbblue}
\setbeamercolor{itemize subitem}{fg=ubbteal}
\setbeamercolor{enumerate item}{fg=ubbblue}

\setbeamerfont{headline}{size=\scriptsize}
\setbeamerfont{footline}{size=\scriptsize}
\setbeamerfont{block title}{size=\normalsize}

% Badan blok rata kiri (bukan rata kanan-kiri): teks slide lebih mudah dibaca
% dan baris terakhir yang pendek tidak lagi memicu "Underfull \hbox".
\addtobeamertemplate{block body begin}{\raggedright}{}
\addtobeamertemplate{block alerted body begin}{\raggedright}{}
\addtobeamertemplate{block example body begin}{\raggedright}{}

% Header: nama mata kuliah dan kode MK
\setbeamertemplate{headline}{%
    \begin{beamercolorbox}[wd=\paperwidth,ht=2.6ex,dp=1.2ex,leftskip=1em,rightskip=1em]{palette primary}%
        \usebeamerfont{headline}\CourseName\ \textperiodcentered\ \CourseCode%
        \hfill \CourseProgram%
    \end{beamercolorbox}%
}

% Footer: minggu dan nomor halaman
\setbeamertemplate{footline}{%
    \begin{beamercolorbox}[wd=\paperwidth,ht=2.4ex,dp=1.2ex,leftskip=1em,rightskip=1em]{palette secondary}%
        \usebeamerfont{footline}\WeekLabel\ \textperiodcentered\ \WeekTopic%
        \hfill \insertframenumber\,/\,\inserttotalframenumber%
    \end{beamercolorbox}%
}

% Halaman judul: judul, subjudul, lalu blok identitas mata kuliah
\setbeamertemplate{title page}{%
    \begin{centering}
        \vskip1.0em
        {\usebeamerfont{title}\usebeamercolor[fg]{title}\inserttitle\par}
        \vskip0.3em
        {\usebeamerfont{subtitle}\usebeamercolor[fg]{subtitle}\insertsubtitle\par}
        \vskip1.2em
        \textcolor{ubbblue}{\rule{0.72\textwidth}{0.8pt}}
        \vskip1.2em
        {\small
            \textbf{\CourseName}\ (\CourseCode)\par
            \vskip0.4em
            \CourseProgram\par
            \CourseFaculty\par
            \CourseInstitution\par
            \vskip0.6em
            \CourseSemester\ \textperiodcentered\ \WeekLabel\par
            \vskip0.8em
            Dosen Pengampu: \CourseLecturer\par
        }
        \vskip1.0em
        \textcolor{ubbblue}{\rule{0.72\textwidth}{0.4pt}}
    \end{centering}
}

% =====================================================================
%  PERINTAH BANTU
% =====================================================================
\newcommand{\R}{\mathbb{R}}
\newcommand{\N}{\mathbb{N}}
\newcommand{\Z}{\mathbb{Z}}
\newcommand{\code}[1]{\texttt{#1}}

% Rujukan gambar dari buku: \bookfig[lebar]{NN}{namaberkas.png}
% Jalur relatif terhadap folder ppt/ (dikompilasi dari sana).
\newcommand{\bookfig}[3][0.80\linewidth]{%
    \includegraphics[width=#1]{../buku_ajar/chapters/#2/figures/#3}%
}

% Keterangan sumber gambar, sepadan dengan \sumber di buku
\newcommand{\figsumber}[1]{%
    \par\vspace{1pt}{\tiny\itshape Sumber: #1\par}%
}

% Penanda blok materi yang sering dipakai di slide
\newcommand{\slidejudul}[1]{{\large\bfseries\textcolor{ubbblue}{#1}}\par\vspace{2pt}}
\newcommand{\kunci}[1]{\textcolor{ubbred}{\textbf{#1}}}

%      \title{...} \subtitle{...} \setWeekInfo{Minggu n}{Topik}
%
%  Berkas minggu hanya berisi \setWeekInfo dan lingkungan frame.
%  Kompilasi dijalankan DARI DALAM folder ppt/ (lihat compile-ppt.bat),
%  karena gambar buku dirujuk dengan jalur relatif ../buku_ajar/.
% =====================================================================

% ---------------------------------------------------------------- bahasa
\usepackage[utf8]{inputenc}
\usepackage[T1]{fontenc}
\usepackage[indonesian]{babel}
\usepackage{lmodern}
\usepackage{microtype}

% ------------------------------------------------------------- matematika
\usepackage{amsmath,amssymb,mathtools}

% ----------------------------------------------------------------- tabel
\usepackage{booktabs}
\usepackage{array}
\usepackage{tabularx}

% --------------------------------------------------------------- grafik
\usepackage{tikz}
\usetikzlibrary{arrows.meta,positioning,shapes.geometric,calc,fit,backgrounds}

% ---------------------------------------------------------- kode program
\usepackage{listings}
\lstdefinestyle{slidecode}{
    basicstyle=\ttfamily\scriptsize,
    keywordstyle=\color{ubbblue}\bfseries,
    commentstyle=\color{ubbgreen}\itshape,
    stringstyle=\color{ubbred},
    numbers=left,
    numberstyle=\tiny\color{gray},
    numbersep=6pt,
    frame=single,
    rulecolor=\color{ubbblue!40},
    backgroundcolor=\color{ubbblue!3},
    breaklines=true,
    showstringspaces=false,
    tabsize=4,
    columns=fullflexible,
    keepspaces=true,
}
\lstset{style=slidecode}

% =====================================================================
%  WARNA INSTITUSI
% =====================================================================
\definecolor{ubbblue}{RGB}{12, 62, 120}      % warna utama
\definecolor{ubbteal}{RGB}{0, 118, 110}      % aksen kedua
\definecolor{ubbgreen}{RGB}{22, 122, 74}     % untuk blok contoh
\definecolor{ubbred}{RGB}{168, 42, 42}       % untuk blok peringatan
\definecolor{ubbgray}{RGB}{70, 76, 84}       % teks sekunder

% =====================================================================
%  IDENTITAS MATA KULIAH
% =====================================================================
\newcommand{\CourseName}{Pengantar Kriptografi}
\newcommand{\CourseCode}{MPFA014}
\newcommand{\CourseSemester}{Semester Ganjil}
\newcommand{\CourseProgram}{Program Studi Matematika}
\newcommand{\CourseFaculty}{Fakultas Matematika dan Ilmu Pengetahuan Alam}
\newcommand{\CourseInstitution}{Universitas Bale Bandung}
\newcommand{\CourseLecturer}{Cecep Suwanda, S.Si., M.Kom.}

% Label minggu, disetel tiap berkas lewat \setWeekInfo
\newcommand{\WeekLabel}{}
\newcommand{\WeekTopic}{}
\newcommand{\setWeekInfo}[2]{%
    \renewcommand{\WeekLabel}{#1}%
    \renewcommand{\WeekTopic}{#2}%
}

% =====================================================================
%  TEMA: header, footer, warna blok
% =====================================================================
\usetheme{default}
\setbeamertemplate{navigation symbols}{}
\setbeamertemplate{blocks}[rounded][shadow=false]

\setbeamercolor{palette primary}{bg=ubbblue, fg=white}
\setbeamercolor{palette secondary}{bg=ubbgray, fg=white}
\setbeamercolor{palette tertiary}{bg=ubbteal, fg=white}

\setbeamercolor{structure}{fg=ubbblue}
\setbeamercolor{normal text}{fg=black!92}
\setbeamercolor{alerted text}{fg=ubbred}

\setbeamercolor{block title}{bg=ubbblue, fg=white}
\setbeamercolor{block body}{bg=ubbblue!6, fg=black!92}
\setbeamercolor{block title alerted}{bg=ubbred, fg=white}
\setbeamercolor{block body alerted}{bg=ubbred!6, fg=black!92}
\setbeamercolor{block title example}{bg=ubbgreen, fg=white}
\setbeamercolor{block body example}{bg=ubbgreen!6, fg=black!92}

\setbeamercolor{itemize item}{fg=ubbblue}
\setbeamercolor{itemize subitem}{fg=ubbteal}
\setbeamercolor{enumerate item}{fg=ubbblue}

\setbeamerfont{headline}{size=\scriptsize}
\setbeamerfont{footline}{size=\scriptsize}
\setbeamerfont{block title}{size=\normalsize}

% Badan blok rata kiri (bukan rata kanan-kiri): teks slide lebih mudah dibaca
% dan baris terakhir yang pendek tidak lagi memicu "Underfull \hbox".
\addtobeamertemplate{block body begin}{\raggedright}{}
\addtobeamertemplate{block alerted body begin}{\raggedright}{}
\addtobeamertemplate{block example body begin}{\raggedright}{}

% Header: nama mata kuliah dan kode MK
\setbeamertemplate{headline}{%
    \begin{beamercolorbox}[wd=\paperwidth,ht=2.6ex,dp=1.2ex,leftskip=1em,rightskip=1em]{palette primary}%
        \usebeamerfont{headline}\CourseName\ \textperiodcentered\ \CourseCode%
        \hfill \CourseProgram%
    \end{beamercolorbox}%
}

% Footer: minggu dan nomor halaman
\setbeamertemplate{footline}{%
    \begin{beamercolorbox}[wd=\paperwidth,ht=2.4ex,dp=1.2ex,leftskip=1em,rightskip=1em]{palette secondary}%
        \usebeamerfont{footline}\WeekLabel\ \textperiodcentered\ \WeekTopic%
        \hfill \insertframenumber\,/\,\inserttotalframenumber%
    \end{beamercolorbox}%
}

% Halaman judul: judul, subjudul, lalu blok identitas mata kuliah
\setbeamertemplate{title page}{%
    \begin{centering}
        \vskip1.0em
        {\usebeamerfont{title}\usebeamercolor[fg]{title}\inserttitle\par}
        \vskip0.3em
        {\usebeamerfont{subtitle}\usebeamercolor[fg]{subtitle}\insertsubtitle\par}
        \vskip1.2em
        \textcolor{ubbblue}{\rule{0.72\textwidth}{0.8pt}}
        \vskip1.2em
        {\small
            \textbf{\CourseName}\ (\CourseCode)\par
            \vskip0.4em
            \CourseProgram\par
            \CourseFaculty\par
            \CourseInstitution\par
            \vskip0.6em
            \CourseSemester\ \textperiodcentered\ \WeekLabel\par
            \vskip0.8em
            Dosen Pengampu: \CourseLecturer\par
        }
        \vskip1.0em
        \textcolor{ubbblue}{\rule{0.72\textwidth}{0.4pt}}
    \end{centering}
}

% =====================================================================
%  PERINTAH BANTU
% =====================================================================
\newcommand{\R}{\mathbb{R}}
\newcommand{\N}{\mathbb{N}}
\newcommand{\Z}{\mathbb{Z}}
\newcommand{\code}[1]{\texttt{#1}}

% Rujukan gambar dari buku: \bookfig[lebar]{NN}{namaberkas.png}
% Jalur relatif terhadap folder ppt/ (dikompilasi dari sana).
\newcommand{\bookfig}[3][0.80\linewidth]{%
    \includegraphics[width=#1]{../buku_ajar/chapters/#2/figures/#3}%
}

% Keterangan sumber gambar, sepadan dengan \sumber di buku
\newcommand{\figsumber}[1]{%
    \par\vspace{1pt}{\tiny\itshape Sumber: #1\par}%
}

% Penanda blok materi yang sering dipakai di slide
\newcommand{\slidejudul}[1]{{\large\bfseries\textcolor{ubbblue}{#1}}\par\vspace{2pt}}
\newcommand{\kunci}[1]{\textcolor{ubbred}{\textbf{#1}}}

%      \title{...} \subtitle{...} \setWeekInfo{Minggu n}{Topik}
%
%  Berkas minggu hanya berisi \setWeekInfo dan lingkungan frame.
%  Kompilasi dijalankan DARI DALAM folder ppt/ (lihat compile-ppt.bat),
%  karena gambar buku dirujuk dengan jalur relatif ../buku_ajar/.
% =====================================================================

% ---------------------------------------------------------------- bahasa
\usepackage[utf8]{inputenc}
\usepackage[T1]{fontenc}
\usepackage[indonesian]{babel}
\usepackage{lmodern}
\usepackage{microtype}

% ------------------------------------------------------------- matematika
\usepackage{amsmath,amssymb,mathtools}

% ----------------------------------------------------------------- tabel
\usepackage{booktabs}
\usepackage{array}
\usepackage{tabularx}

% --------------------------------------------------------------- grafik
\usepackage{tikz}
\usetikzlibrary{arrows.meta,positioning,shapes.geometric,calc,fit,backgrounds}

% ---------------------------------------------------------- kode program
\usepackage{listings}
\lstdefinestyle{slidecode}{
    basicstyle=\ttfamily\scriptsize,
    keywordstyle=\color{ubbblue}\bfseries,
    commentstyle=\color{ubbgreen}\itshape,
    stringstyle=\color{ubbred},
    numbers=left,
    numberstyle=\tiny\color{gray},
    numbersep=6pt,
    frame=single,
    rulecolor=\color{ubbblue!40},
    backgroundcolor=\color{ubbblue!3},
    breaklines=true,
    showstringspaces=false,
    tabsize=4,
    columns=fullflexible,
    keepspaces=true,
}
\lstset{style=slidecode}

% =====================================================================
%  WARNA INSTITUSI
% =====================================================================
\definecolor{ubbblue}{RGB}{12, 62, 120}      % warna utama
\definecolor{ubbteal}{RGB}{0, 118, 110}      % aksen kedua
\definecolor{ubbgreen}{RGB}{22, 122, 74}     % untuk blok contoh
\definecolor{ubbred}{RGB}{168, 42, 42}       % untuk blok peringatan
\definecolor{ubbgray}{RGB}{70, 76, 84}       % teks sekunder

% =====================================================================
%  IDENTITAS MATA KULIAH
% =====================================================================
\newcommand{\CourseName}{Pengantar Kriptografi}
\newcommand{\CourseCode}{MPFA014}
\newcommand{\CourseSemester}{Semester Ganjil}
\newcommand{\CourseProgram}{Program Studi Matematika}
\newcommand{\CourseFaculty}{Fakultas Matematika dan Ilmu Pengetahuan Alam}
\newcommand{\CourseInstitution}{Universitas Bale Bandung}
\newcommand{\CourseLecturer}{Cecep Suwanda, S.Si., M.Kom.}

% Label minggu, disetel tiap berkas lewat \setWeekInfo
\newcommand{\WeekLabel}{}
\newcommand{\WeekTopic}{}
\newcommand{\setWeekInfo}[2]{%
    \renewcommand{\WeekLabel}{#1}%
    \renewcommand{\WeekTopic}{#2}%
}

% =====================================================================
%  TEMA: header, footer, warna blok
% =====================================================================
\usetheme{default}
\setbeamertemplate{navigation symbols}{}
\setbeamertemplate{blocks}[rounded][shadow=false]

\setbeamercolor{palette primary}{bg=ubbblue, fg=white}
\setbeamercolor{palette secondary}{bg=ubbgray, fg=white}
\setbeamercolor{palette tertiary}{bg=ubbteal, fg=white}

\setbeamercolor{structure}{fg=ubbblue}
\setbeamercolor{normal text}{fg=black!92}
\setbeamercolor{alerted text}{fg=ubbred}

\setbeamercolor{block title}{bg=ubbblue, fg=white}
\setbeamercolor{block body}{bg=ubbblue!6, fg=black!92}
\setbeamercolor{block title alerted}{bg=ubbred, fg=white}
\setbeamercolor{block body alerted}{bg=ubbred!6, fg=black!92}
\setbeamercolor{block title example}{bg=ubbgreen, fg=white}
\setbeamercolor{block body example}{bg=ubbgreen!6, fg=black!92}

\setbeamercolor{itemize item}{fg=ubbblue}
\setbeamercolor{itemize subitem}{fg=ubbteal}
\setbeamercolor{enumerate item}{fg=ubbblue}

\setbeamerfont{headline}{size=\scriptsize}
\setbeamerfont{footline}{size=\scriptsize}
\setbeamerfont{block title}{size=\normalsize}

% Badan blok rata kiri (bukan rata kanan-kiri): teks slide lebih mudah dibaca
% dan baris terakhir yang pendek tidak lagi memicu "Underfull \hbox".
\addtobeamertemplate{block body begin}{\raggedright}{}
\addtobeamertemplate{block alerted body begin}{\raggedright}{}
\addtobeamertemplate{block example body begin}{\raggedright}{}

% Header: nama mata kuliah dan kode MK
\setbeamertemplate{headline}{%
    \begin{beamercolorbox}[wd=\paperwidth,ht=2.6ex,dp=1.2ex,leftskip=1em,rightskip=1em]{palette primary}%
        \usebeamerfont{headline}\CourseName\ \textperiodcentered\ \CourseCode%
        \hfill \CourseProgram%
    \end{beamercolorbox}%
}

% Footer: minggu dan nomor halaman
\setbeamertemplate{footline}{%
    \begin{beamercolorbox}[wd=\paperwidth,ht=2.4ex,dp=1.2ex,leftskip=1em,rightskip=1em]{palette secondary}%
        \usebeamerfont{footline}\WeekLabel\ \textperiodcentered\ \WeekTopic%
        \hfill \insertframenumber\,/\,\inserttotalframenumber%
    \end{beamercolorbox}%
}

% Halaman judul: judul, subjudul, lalu blok identitas mata kuliah
\setbeamertemplate{title page}{%
    \begin{centering}
        \vskip1.0em
        {\usebeamerfont{title}\usebeamercolor[fg]{title}\inserttitle\par}
        \vskip0.3em
        {\usebeamerfont{subtitle}\usebeamercolor[fg]{subtitle}\insertsubtitle\par}
        \vskip1.2em
        \textcolor{ubbblue}{\rule{0.72\textwidth}{0.8pt}}
        \vskip1.2em
        {\small
            \textbf{\CourseName}\ (\CourseCode)\par
            \vskip0.4em
            \CourseProgram\par
            \CourseFaculty\par
            \CourseInstitution\par
            \vskip0.6em
            \CourseSemester\ \textperiodcentered\ \WeekLabel\par
            \vskip0.8em
            Dosen Pengampu: \CourseLecturer\par
        }
        \vskip1.0em
        \textcolor{ubbblue}{\rule{0.72\textwidth}{0.4pt}}
    \end{centering}
}

% =====================================================================
%  PERINTAH BANTU
% =====================================================================
\newcommand{\R}{\mathbb{R}}
\newcommand{\N}{\mathbb{N}}
\newcommand{\Z}{\mathbb{Z}}
\newcommand{\code}[1]{\texttt{#1}}

% Rujukan gambar dari buku: \bookfig[lebar]{NN}{namaberkas.png}
% Jalur relatif terhadap folder ppt/ (dikompilasi dari sana).
\newcommand{\bookfig}[3][0.80\linewidth]{%
    \includegraphics[width=#1]{../buku_ajar/chapters/#2/figures/#3}%
}

% Keterangan sumber gambar, sepadan dengan \sumber di buku
\newcommand{\figsumber}[1]{%
    \par\vspace{1pt}{\tiny\itshape Sumber: #1\par}%
}

% Penanda blok materi yang sering dipakai di slide
\newcommand{\slidejudul}[1]{{\large\bfseries\textcolor{ubbblue}{#1}}\par\vspace{2pt}}
\newcommand{\kunci}[1]{\textcolor{ubbred}{\textbf{#1}}}


\title{Menghitung dengan Titik, Bukan Bilangan Besar}
\subtitle{Kurva eliptik, ECDLP, ECDH, dan ECEG}
\setWeekInfo{Minggu 13}{Kriptografi kurva eliptik}

\begin{document}

% ----------------------------------------------------------- halaman judul
\begin{frame}[plain]
    \titlepage
\end{frame}

% ------------------------------------------------------------------ peta
\begin{frame}{Peta Materi Pertemuan Ini}
    \begin{columns}[T,onlytextwidth]
        \begin{column}{0.52\textwidth}
            \slidejudul{Lima bagian}
            \footnotesize
            \begin{enumerate}
                \item \textbf{Landasan yang dipakai} --- grup, medan
                      berhingga, dan invers modulo.
                \item \textbf{Kurva eliptik} --- bentuk kurva, titik
                      $\mathcal{O}$, penjumlahan, dan penggandaan titik.
                \item \textbf{ECDLP} --- persoalan matematis yang menopang
                      keamanan ECC.
                \item \textbf{ECDH dan ECEG} --- kesepakatan kunci dan
                      enkripsi dengan titik kurva.
                \item \textbf{Keamanan dan keuntungan} --- ukuran kunci serta
                      pemakaian ECC.
            \end{enumerate}
        \end{column}
        \begin{column}{0.44\textwidth}
            \begin{block}{Referensi}
                \footnotesize
                Bab 14 buku ajar; Munir, \textit{IF4020 Kriptografi}, ITB,
                2026 (modul 17 dan 18)
            \end{block}
            \begin{exampleblock}{Benang merah}
                \footnotesize
                ECC tidak memperbesar bilangan. Ia memindahkan hitungan ke
                \kunci{titik-titik pada kurva berhingga}.
            \end{exampleblock}
        \end{column}
    \end{columns}
\end{frame}

% ---------------------------------------------------- capaian pembelajaran
\begin{frame}{Capaian Pembelajaran}
    \begin{block}{Sub-CPMK yang dibebankan}
        \footnotesize
        \textbf{Sub-CPMK 1.4} --- Mengenal algoritma kriptografi kunci
        publik, yaitu RSA, Diffie--Hellman, ElGamal, knapsack, dan
        \textit{Elliptic Curve Cryptography} (ECC), beserta persoalan
        matematis yang mendasarinya.
    \end{block}
    \vspace{0.15cm}
    \begin{exampleblock}{Setelah pertemuan ini, mahasiswa mampu}
        \footnotesize
        \begin{itemize}
            \item menjelaskan aksioma grup dan medan berhingga sejauh dipakai
                  oleh operasi titik;
            \item menjelaskan aturan penjumlahan dan penggandaan titik pada
                  kurva eliptik;
            \item menjelaskan ECDLP sebagai persoalan yang mendasari keamanan
                  ECC;
            \item menjelaskan langkah ECDH dan ECEG beserta hasil hitungnya;
            \item membandingkan ukuran kunci ECC dengan RSA pada tingkat
                  keamanan yang setara.
        \end{itemize}
    \end{exampleblock}
\end{frame}

% =====================================================================
\section{Landasan yang Dipakai}
% =====================================================================

\begin{frame}{Mengapa Hitungannya Dipindahkan ke Titik}
    \begin{columns}[T,onlytextwidth]
        \begin{column}{0.54\textwidth}
            \footnotesize
            \begin{itemize}
                \item Sistem kunci publik biasa menghitung di dalam himpunan
                      bilangan modulo.
                \item ECC menghitung dengan pasangan $(x,y)$ yang memenuhi
                      persamaan sebuah kurva.
                \item Pasangan itu harus dapat dijumlahkan dan diulang
                      berkali-kali.
                \item Karena itu himpunan titiknya harus berbentuk
                      \kunci{grup}.
            \end{itemize}
        \end{column}
        \begin{column}{0.42\textwidth}
            \begin{block}{Syarat pendukung}
                \footnotesize
                Koordinat titik diambil dari medan berhingga. Setiap operasi
                memberi hasil eksak, tanpa pembulatan.
            \end{block}
            \begin{alertblock}{Satu operasi dasar}
                \footnotesize
                Seluruh protokol ECC dibangun dari perkalian titik $kP$.
            \end{alertblock}
        \end{column}
    \end{columns}
    \vspace{0.1cm}
    \begin{exampleblock}{Peran bilangan riil}
        \footnotesize
        Kurva di atas $\R$ dipakai untuk memahami geometri. Kriptografi
        memakai kurva di atas medan berhingga, sebab bilangan riil dibulatkan
        dan galatnya menumpuk.
    \end{exampleblock}
\end{frame}

\begin{frame}{Grup: Empat Aksioma}
    \begin{block}{Definisi}
        \footnotesize
        Grup adalah himpunan $G$ dengan operasi biner $*$ yang memetakan
        setiap pasangan elemen $G$ kembali ke $G$.
    \end{block}
    \vspace{0.05cm}
    \begin{enumerate}
        \item \textbf{Ketertutupan}: $a * b \in G$ untuk semua $a,b \in G$.
        \item \textbf{Asosiatif}: $a * (b * c) = (a * b) * c$.
        \item \textbf{Elemen netral}: ada $e \in G$ dengan
              $a * e = e * a = a$.
        \item \textbf{Elemen invers}: setiap $a$ punya $a'$ dengan
              $a * a' = e$.
    \end{enumerate}
    \vspace{0.05cm}
    \begin{exampleblock}{Contoh: $\langle \Z_{11}, + \rangle$}
        \footnotesize
        Elemen netralnya $0$. Invers dari $a$ adalah $11-a$, sehingga
        $7 + 4 = 11 \equiv 0 \pmod{11}$.
    \end{exampleblock}
\end{frame}

\begin{frame}{Grup Abelian dan Satu Sifat yang Dipakai ECDH}
    \begin{block}{Aksioma kelima}
        \footnotesize
        Grup disebut \textbf{abelian} bila operasinya komutatif:
        $a * b = b * a$ untuk semua $a,b \in G$.
    \end{block}
    \vspace{0.1cm}
    \begin{itemize}
        \item Penjumlahan modulo $n$ bersifat komutatif.
        \item Perkalian modulo prima juga bersifat komutatif.
        \item Urutan pengoperasian karena itu tidak mengubah hasil.
        \item ECDH memerlukan sifat ini untuk menyamakan
              $a(bB)$ dengan $b(aB)$.
        \item Tanpa komutatif, kedua pihak tidak akan bertemu di titik yang
              sama.
    \end{itemize}
    \[
        a(bB) = (ab)B = (ba)B = b(aB) .
    \]
\end{frame}

\begin{frame}{Medan Berhingga $GF(p)$ dan Invers Modulo}
    \begin{block}{Definisi medan}
        \footnotesize
        Himpunan $F$ dengan dua operasi $+$ dan $\cdot$ disebut medan bila
        $\langle F,+ \rangle$ dan $\langle F-\{0\},\cdot \rangle$ keduanya
        grup abelian, serta perkalian menyebar terhadap penjumlahan.
    \end{block}
    \vspace{0.05cm}
    \begin{itemize}
        \item $GF(p)$ berisi $\{0,1,\dots,p-1\}$ dengan operasi modulo prima
              $p$.
        \item Setiap unsur bukan nol punya invers:
              $a \cdot a^{-1} \equiv 1 \pmod p$.
        \item Pembagian pada rumus titik diganti perkalian dengan invers.
        \item Contoh: $3^{-1} \equiv 4 \pmod{11}$ sebab
              $3 \cdot 4 = 12 \equiv 1$.
    \end{itemize}
    \begin{alertblock}{Medan biner}
        \footnotesize
        ECC juga dapat didefinisikan di $GF(2^m)$. Pertemuan ini memakai
        $GF(p)$ agar operasinya tetap aritmetika modulo biasa.
    \end{alertblock}
\end{frame}

% =====================================================================
\section{Kurva Eliptik}
% =====================================================================

\begin{frame}{Bentuk Kurva dan Syaratnya}
    \begin{block}{Persamaan kurva}
        \[
            y^2 = x^3 + ax + b,
            \qquad 4a^3 + 27b^2 \neq 0 .
        \]
    \end{block}
    \vspace{0.05cm}
    \begin{itemize}
        \item Nilai $a$ dan $b$ menentukan bentuk serta letak kurva.
        \item Syarat diskriminan menolak kurva yang memiliki akar ganda.
        \item Contoh: $y^2 = x^3 - 4x$ memakai $a=-4$ dan $b=0$.
        \item Contoh: $y^2 = x^3 - 12x + 13$ memakai $a=-12$ dan $b=13$.
    \end{itemize}
    \begin{alertblock}{Kurva di atas $GF(p)$}
        \footnotesize
        Untuk kriptografi persamaannya dibaca modulo $p$:
        $y^2 \equiv x^3 + ax + b \pmod p$, dengan
        $x,y \in \{0,1,\dots,p-1\}$.
    \end{alertblock}
\end{frame}

\begin{frame}{Lima Bentuk Kurva di Atas Bilangan Riil}
    \centering
    \bookfig[0.94\linewidth]{14}{galeri-kurva.png}
    \figsumber{Munir, \textit{IF4020 Kriptografi}, ITB.}
    \vspace{0.1cm}
    \begin{itemize}
        \footnotesize
        \item Setiap pasangan $(a,b)$ memberi bentuk yang berbeda.
        \item Jumlah komponen kurva dapat berubah dari satu pasangan ke
              pasangan lain.
        \item Gambar ini menjelaskan geometri, bukan dipakai menghitung.
    \end{itemize}
\end{frame}

\begin{frame}{Titik di Ketakhinggaan sebagai Elemen Netral}
    \begin{columns}[T,onlytextwidth]
        \begin{column}{0.54\textwidth}
            \footnotesize
            \begin{itemize}
                \item Titik-titik biasa belum menyediakan elemen netral.
                \item Ditambahkan satu titik $\mathcal{O}$, yaitu titik di
                      ketakhinggaan.
                \item $\mathcal{O}$ tidak tergambar, tetapi melengkapi
                      grup.
                \item Lawan titik diperoleh dengan pencerminan:
                      $-P = (x,(-y) \bmod p)$.
            \end{itemize}
        \end{column}
        \begin{column}{0.42\textwidth}
            \begin{block}{Peran $\mathcal{O}$}
                \[
                    P + \mathcal{O} = \mathcal{O} + P = P
                \]
                \[
                    P + (-P) = \mathcal{O}
                \]
            \end{block}
        \end{column}
    \end{columns}
    \vspace{0.1cm}
    \begin{alertblock}{Akibatnya}
        \footnotesize
        Penjumlahan dua titik yang ordinatnya berlawanan tidak menghasilkan
        titik baru, melainkan $\mathcal{O}$.
    \end{alertblock}
\end{frame}

\begin{frame}{Penjumlahan Titik: Empat Langkah Geometris}
    \begin{enumerate}
        \item Tarik garis lurus melalui $P$ dan $Q$.
        \item Garis itu memotong kurva di titik ketiga, yaitu $-R$.
        \item Cerminkan $-R$ terhadap sumbu-$x$.
        \item Hasil pencerminan adalah $R = P + Q$.
    \end{enumerate}
    \vspace{0.15cm}
    \begin{columns}[T,onlytextwidth]
        \begin{column}{0.54\textwidth}
            \begin{block}{Koordinat}
                Bila $R=(x_r,y_r)$, maka
                \[
                    -R = (x_r,-y_r) \pmod p .
                \]
            \end{block}
        \end{column}
        \begin{column}{0.42\textwidth}
            \begin{alertblock}{Inti langkah}
                \footnotesize
                Pencerminan inilah yang membedakan penjumlahan titik dari
                penjumlahan bilangan biasa.
            \end{alertblock}
        \end{column}
    \end{columns}
\end{frame}

\begin{frame}{Dari Geometri ke Rumus: Hubungan Vieta}
    \footnotesize
    \begin{itemize}
        \item Garis melalui $P(x_p,y_p)$ dan $Q(x_q,y_q)$ ditulis
              $y = mx + c$.
        \item Substitusi ke persamaan kurva menghasilkan polinom kubik
              \[
                  x^3 - m^2x^2 + (a-2mc)x + (b-c^2) = 0 .
              \]
        \item Akar-akarnya adalah $x_p$, $x_q$, dan $x_r$.
        \item Vieta memberi $x_p + x_q + x_r = m^2$.
        \item Karena itu
              $x_r = m^2 - x_p - x_q$.
        \item Ordinatnya $y_r = m(x_p - x_r) - y_p$.
        \item Di atas $GF(p)$, setiap hasil diambil modulo $p$.
    \end{itemize}
    \begin{exampleblock}{Gradien garis potong}
        \[
            m = \frac{y_p - y_q}{x_p - x_q}
              \equiv (y_p - y_q)(x_p - x_q)^{-1} \pmod p .
        \]
    \end{exampleblock}
\end{frame}

\begin{frame}{Rumus Penjumlahan di $GF(p)$}
    \begin{block}{Untuk $P \neq Q$}
        \[
        \begin{aligned}
            m   &\equiv (y_p - y_q)(x_p - x_q)^{-1} \pmod p, \\
            x_r &\equiv m^2 - x_p - x_q \pmod p, \\
            y_r &\equiv m(x_p - x_r) - y_p \pmod p .
        \end{aligned}
        \]
    \end{block}
    \vspace{0.1cm}
    \begin{itemize}
        \item Pembagian diganti perkalian dengan invers modulo.
        \item Rumus hanya memakai operasi yang tersedia di dalam medan.
        \item Hasilnya selalu titik lain pada kurva yang sama.
    \end{itemize}
    \begin{alertblock}{Kasus khusus}
        \footnotesize
        Bila $x_p = x_q$ dan $y_p = -y_q$, garisnya tegak dan
        $P + Q = \mathcal{O}$.
    \end{alertblock}
\end{frame}

\begin{frame}{Penggandaan Titik}
    \begin{block}{Menjumlahkan titik dengan dirinya sendiri}
        Garis yang dipakai adalah garis singgung di $P$, sehingga
        \[
            m \equiv (3x_p^2 + a)(2y_p)^{-1} \pmod p .
        \]
    \end{block}
    \vspace{0.05cm}
    \begin{block}{Koordinat hasil $2P$}
        \[
            x_r \equiv m^2 - 2x_p \pmod p,
            \qquad
            y_r \equiv m(x_p - x_r) - y_p \pmod p .
        \]
    \end{block}
    \vspace{0.05cm}
    \begin{alertblock}{Bila $y_p = 0$}
        \footnotesize
        Penyebut gradien menjadi nol. Garis singgungnya tegak, sehingga
        $2P = \mathcal{O}$.
    \end{alertblock}
\end{frame}

\begin{frame}{Perkalian Skalar dengan Penjumlahan dan Penggandaan}
    \begin{block}{Pelelaran titik}
        \[
            kP = \underbrace{P + P + \dots + P}_{k \text{ kali}} .
        \]
    \end{block}
    \vspace{0.05cm}
    \begin{itemize}
        \item Penjumlahan dan penggandaan adalah dua operasi dasar.
        \item Keduanya dipakai bergantian mengikuti bit skalar $k$.
        \item Contoh: $23P$ dihitung tanpa menjumlahkan $P$ dua puluh dua
              kali.
    \end{itemize}
    \[
        23P = 2\bigl(2(2(2P) + P) + P\bigr) + P .
    \]
    \begin{exampleblock}{Mengapa penting}
        \footnotesize
        ECDH dan ECEG sama-sama memerlukan $kP$. Biaya hitungnya tumbuh
        sebanding dengan panjang $k$ dalam bit, bukan dengan nilai $k$.
    \end{exampleblock}
\end{frame}

\begin{frame}{Titik Kurva Membentuk Grup Abelian}
    \footnotesize
    \begin{tabularx}{\linewidth}{@{}l l >{\raggedright\arraybackslash}X@{}}
        \toprule
        \textbf{Aksioma} & \textbf{Pemenuhan} & \textbf{Alasan} \\
        \midrule
        Ketertutupan & hasil $P+Q$ &
        rumus koordinat hanya memakai operasi medan \\
        \addlinespace
        Asosiatif & --- &
        $P+(Q+R) = (P+Q)+R$ \\
        \addlinespace
        Elemen netral & $\mathcal{O}$ &
        $P+\mathcal{O} = \mathcal{O}+P = P$ \\
        \addlinespace
        Elemen invers & $-P=(x_p,-y_p)$ &
        $P+(-P) = \mathcal{O}$ \\
        \addlinespace
        Komutatif & --- &
        $P+Q = Q+P$ \\
        \bottomrule
    \end{tabularx}
    \vspace{0.25cm}
    \begin{exampleblock}{Kesimpulan}
        \footnotesize
        Himpunan semua titik kurva ditambah $\mathcal{O}$, dengan operasi
        penjumlahan titik, adalah \kunci{grup abelian}.
    \end{exampleblock}
\end{frame}

\begin{frame}{Menghitung Titik pada Kurva di $GF(11)$}
    \begin{columns}[T,onlytextwidth]
        \begin{column}{0.56\textwidth}
            \footnotesize
            Kurva: $y^2 \equiv x^3 + x + 6 \pmod{11}$.
            \begin{itemize}
                \item Untuk tiap $x$, hitung ruas kanan modulo $11$.
                \item Nilai $y^2$ yang mungkin: $0,1,3,4,5,9$.
                \item Hanya $x=2,3,5,7,8,10$ yang memberi kuadrat.
                \item Setiap $x$ itu memberi dua nilai $y$, kecuali bila
                      ruas kanannya nol.
            \end{itemize}
            \begin{exampleblock}{Hasil}
                \footnotesize
                Ada $12$ titik berhingga. Dengan $\mathcal{O}$, orde grupnya
                $13$.
            \end{exampleblock}
        \end{column}
        \begin{column}{0.40\textwidth}
            \centering
            \bookfig[0.98\linewidth]{14}{sebaran-titik-gf11.png}
            \figsumber{Munir, \textit{IF4020 Kriptografi}, ITB.}
        \end{column}
    \end{columns}
\end{frame}

% =====================================================================
\section{ECDLP dan Pembangkitan Kunci}
% =====================================================================

\begin{frame}{ECDLP: Membalik Perkalian Titik}
    \begin{block}{Persoalan logaritma diskret kurva eliptik}
        \footnotesize
        Diberikan titik $P$ dan $Q$ pada sebuah kurva. Carilah bilangan bulat
        $k$ sehingga
        \[
            Q = kP .
        \]
    \end{block}
    \vspace{0.1cm}
    \begin{columns}[T,onlytextwidth]
        \begin{column}{0.48\textwidth}
            \begin{exampleblock}{Arah maju: mudah}
                \footnotesize
                Dari $k$ dan $P$, titik $Q$ dihitung dengan penjumlahan dan
                penggandaan.
            \end{exampleblock}
        \end{column}
        \begin{column}{0.48\textwidth}
            \begin{alertblock}{Arah balik: sukar}
                \footnotesize
                Dari $P$ dan $Q$, mencari $k$ sangat sukar bila orde grup
                besar.
            \end{alertblock}
        \end{column}
    \end{columns}
    \vspace{0.1cm}
    \begin{itemize}
        \item Bilangan $k$ disebut logaritma diskret dari $Q$ dengan basis
              $P$.
        \item Keamanan ECC bersandar pada kesukaran arah balik ini.
    \end{itemize}
\end{frame}

\begin{frame}{Membangkitkan Pasangan Kunci ECC}
    \begin{enumerate}
        \item Sepakati kurva $y^2 \equiv x^3 + ax + b \pmod p$, yaitu nilai
              $a$, $b$, dan $p$.
        \item Sepakati grup titik yang dihasilkan kurva itu.
        \item Pilih titik basis $B(x_B,y_B)$ dari grup tersebut.
    \end{enumerate}
    \vspace{0.1cm}
    \begin{columns}[T,onlytextwidth]
        \begin{column}{0.50\textwidth}
            \begin{block}{Milik pribadi}
                \footnotesize
                Kunci privat adalah bilangan bulat $x$ dengan
                $1 \le x \le p-1$.
            \end{block}
        \end{column}
        \begin{column}{0.46\textwidth}
            \begin{block}{Milik publik}
                \footnotesize
                Kunci publik adalah titik
                \[
                    Q = x \cdot B .
                \]
            \end{block}
        \end{column}
    \end{columns}
    \vspace{0.1cm}
    \begin{alertblock}{Letak kesukarannya}
        \footnotesize
        Siapa pun boleh mengetahui $Q$ dan $B$. Memulihkan $x$ dari keduanya
        berarti memecahkan ECDLP.
    \end{alertblock}
\end{frame}

% =====================================================================
\section{ECDH}
% =====================================================================

\begin{frame}{Pertukaran Kunci ECDH}
    \begin{enumerate}
        \item Alice dan Bob menyepakati kurva serta titik basis $B$.
        \item Alice memilih $a$ dan menghitung $P_A = aB$; Bob memilih $b$
              dan menghitung $P_B = bB$.
        \item Keduanya bertukar $P_A$ dan $P_B$, lalu mengalikan kunci publik
              mitranya dengan kunci privat sendiri.
    \end{enumerate}
    \vspace{0.1cm}
    \begin{block}{Hasil di kedua sisi}
        \[
        \begin{aligned}
            \text{Alice:}\quad K &= a \cdot P_B = a(bB) = abB, \\
            \text{Bob:}\quad   K &= b \cdot P_A = b(aB) = abB .
        \end{aligned}
        \]
    \end{block}
    \vspace{0.05cm}
    \begin{alertblock}{Yang tidak pernah dikirim}
        \footnotesize
        Kunci privat $a$ dan $b$ tidak pernah melewati saluran. Yang berpindah
        hanya titik $P_A$ dan $P_B$.
    \end{alertblock}
\end{frame}

\begin{frame}{Contoh ECDH pada $GF(5)$}
    \begin{columns}[T,onlytextwidth]
        \begin{column}{0.48\textwidth}
            \footnotesize
            Kurva: $y^2 \equiv x^3 + 2x + 1 \pmod 5$.
            \begin{itemize}
                \item Enam titik berhingga, ditambah $\mathcal{O}$ menjadi
                      orde $7$.
                \item Titik basis $B(0,1)$.
                \item Alice memilih $a=2$; Bob memilih $b=3$.
            \end{itemize}
            \vspace{0.1cm}
            \begin{tabular}{@{}c l l@{}}
                \toprule
                $k$ & $kB$ & Peran \\
                \midrule
                $1$ & $(0,1)$ & basis $B$ \\
                $2$ & $(1,3)$ & publik Alice \\
                $3$ & $(3,3)$ & publik Bob \\
                $4$ & $(3,2)$ & --- \\
                $5$ & $(1,2)$ & --- \\
                $6$ & $(0,4)$ & kunci bersama \\
                $7$ & $\mathcal{O}$ & orde tercapai \\
                \bottomrule
            \end{tabular}
        \end{column}
        \begin{column}{0.48\textwidth}
            \begin{block}{Kunci publik}
                \footnotesize
                \[
                    P_A = 2B = (1,3), \qquad
                    P_B = 3B = (3,3) .
                \]
            \end{block}
            \begin{exampleblock}{Kunci bersama}
                \footnotesize
                Alice: $K = 2 \cdot (3,3)$.

                Bob: $K = 3 \cdot (1,3)$.

                Keduanya memperoleh $(0,4) = 6B$. Karena $B$ berorde $7$,
                maka $6B = -B$.
            \end{exampleblock}
        \end{column}
    \end{columns}
\end{frame}

% =====================================================================
\section{ECEG}
% =====================================================================

\begin{frame}{Pesan Harus Menjadi Titik: Metode Koblitz}
    \begin{itemize}
        \item ECEG mengenkripsi titik kurva, bukan rangkaian karakter.
        \item Setiap karakter dipetakan ke bilangan $m$: angka $0$--$9$
              bernilai $0$--$9$, sedangkan huruf \texttt{A}--\texttt{Z}
              bernilai $10$--$35$.
        \item Pilih parameter basis $k$, lalu coba $x = mk+1$, $mk+2$, dan
              seterusnya.
        \item Percobaan berhenti ketika ruas kanan kurva menjadi kuadrat
              modulo $p$.
    \end{itemize}
    \begin{exampleblock}{Huruf \texttt{B} pada $y^2 \equiv x^3 - x + 188
        \pmod{751}$}
        \footnotesize
        Nilai $m=11$ dan $k=20$ memberi $x = 11 \cdot 20 + 1 = 221$.
        Percobaan $x=221,222,223$ gagal. Pada $x=224$ ruas kanan bernilai
        $673$, dan $248^2 \equiv 503^2 \equiv 673$. Jadi huruf \texttt{B}
        menjadi titik $(224,248)$.
    \end{exampleblock}
    \begin{alertblock}{Dekode}
        \footnotesize
        $m = \lfloor (x-1)/k \rfloor = \lfloor 223/20 \rfloor = 11$.
    \end{alertblock}
\end{frame}

\begin{frame}{Enkripsi ECEG}
    \begin{block}{Bahan yang sudah dimiliki Alice}
        \footnotesize
        Kurva, titik basis $B$, dan kunci publik Bob
        $P_B = bB$. Pesan $M$ sudah diubah menjadi titik $P_M$.
    \end{block}
    \vspace{0.1cm}
    \begin{enumerate}
        \item Alice memilih bilangan acak $k$ dengan $1 \le k \le p-1$.
        \item Alice menghitung titik $kB$.
        \item Alice menutupi pesan dengan $P_M + kP_B$.
        \item Cipherteks dikirim sebagai pasangan titik
              \[
                  P_C = \left[\, (kB),\ (P_M + kP_B) \,\right] .
              \]
    \end{enumerate}
    \begin{alertblock}{Peran setiap komponen}
        \footnotesize
        $kB$ adalah bagian sementara yang dibutuhkan penerima. Komponen kedua
        menyembunyikan titik pesan.
    \end{alertblock}
\end{frame}

\begin{frame}{Dekripsi ECEG}
    \begin{block}{Langkah Bob}
        \footnotesize
        Bob memakai kunci privat $b$ pada komponen pertama, lalu
        mengurangkannya dari komponen kedua:
        \[
        \begin{aligned}
            (P_M + kP_B) - b(kB)
                &= P_M + k(bB) - b(kB) \\
                &= P_M + kbB - kbB \\
                &= P_M .
        \end{aligned}
        \]
    \end{block}
    \vspace{0.05cm}
    \begin{itemize}
        \item Pengurangan titik adalah penjumlahan dengan titik lawannya.
        \item Setelah $P_M$ diperoleh, Koblitz dibalik untuk mendapat $M$.
    \end{itemize}
    \begin{exampleblock}{Sumber keberhasilan}
        \footnotesize
        $b(kB) = k(bB)$ berlaku karena perkalian skalar bersifat asosiatif
        dan komutatif.
    \end{exampleblock}
\end{frame}

% =====================================================================
\section{Keamanan dan Keuntungan ECC}
% =====================================================================

\begin{frame}{Keamanan ECC Bersandar pada ECDLP}
    \begin{itemize}
        \item Penyerang melihat kurva, titik basis $B$, dan kunci publik
              $Q = xB$.
        \item Yang dirahasiakan hanyalah skalar $x$.
        \item Untuk memperoleh $x$, penyerang harus menyelesaikan ECDLP.
        \item Kesukarannya bertambah seiring orde titik basis.
    \end{itemize}
    \begin{exampleblock}{Ilustrasi orde}
        \footnotesize
        Pada kurva $GF(11)$ sebelumnya, $P(2,4)$ berorde $13$ sehingga
        $13P = \mathcal{O}$. Kurva yang dipakai sungguhan memilih titik
        basis berorde sangat besar.
    \end{exampleblock}
    \begin{alertblock}{Salah paham yang sering muncul}
        \footnotesize
        Keamanan tidak berasal dari kerahasiaan kurva. Kurva, $p$, dan $B$
        justru parameter publik yang disepakati.
    \end{alertblock}
\end{frame}

\begin{frame}{Ukuran Kunci pada Keamanan Setara}
    \footnotesize
    \begin{tabularx}{\linewidth}{@{}>{\raggedright\arraybackslash}p{5.0cm}
                                  c c >{\raggedright\arraybackslash}X@{}}
        \toprule
        \textbf{Kebutuhan} & \textbf{RSA} & \textbf{ECC} & \textbf{Catatan} \\
        \midrule
        Keamanan setara & 1024 bit & 160 bit &
        kunci ECC jauh lebih pendek \\
        \addlinespace
        Menyalurkan kunci AES-128 & 3072 bit & 256 bit &
        selisihnya makin besar \\
        \bottomrule
    \end{tabularx}
    \vspace{0.3cm}
    \begin{itemize}
        \item Dengan keamanan yang sama, ECC menuntut bit jauh lebih sedikit.
        \item Pertambahan ukuran kunci ECC lebih landai daripada RSA.
        \item Kunci pendek menghemat penyimpanan dan lebar pita.
    \end{itemize}
    \begin{alertblock}{Yang dibandingkan}
        \footnotesize
        Angka di atas adalah ukuran kunci, bukan ukuran cipherteks.
    \end{alertblock}
\end{frame}

\begin{frame}{Keuntungan dan Tempat Pemakaian ECC}
    \begin{columns}[T,onlytextwidth]
        \begin{column}{0.50\textwidth}
            \begin{block}{Keuntungan}
                \footnotesize
                \begin{itemize}
                    \item kunci lebih pendek pada keamanan setara;
                    \item memori dan penyimpanan lebih hemat;
                    \item komputasi lebih ringan;
                    \item lebar pita yang dibutuhkan lebih kecil.
                \end{itemize}
            \end{block}
        \end{column}
        \begin{column}{0.46\textwidth}
            \begin{block}{Pemakaian yang cocok}
                \footnotesize
                \begin{itemize}
                    \item piranti komunikasi nirkabel;
                    \item \textit{smart card};
                    \item peladen web dengan banyak sesi;
                    \item perangkat berdaya dan memori terbatas.
                \end{itemize}
            \end{block}
        \end{column}
    \end{columns}
    \vspace{0.15cm}
    \begin{exampleblock}{Layanan yang diberikan}
        \footnotesize
        Sama seperti sistem kunci publik lain: kerahasiaan, integritas, dan
        autentikasi. Yang berbeda adalah \kunci{biaya} untuk mencapai tingkat
        keamanan itu.
    \end{exampleblock}
\end{frame}

% =====================================================================
\section{Contoh Terhitung}
% =====================================================================

\begin{frame}{Contoh: Penjumlahan $P+Q$ pada $GF(11)$}
    \footnotesize
    Kurva $y^2 \equiv x^3 + x + 6 \pmod{11}$, dengan $a=1$ dan $p=11$.
    Titik $P(2,4)$ dan $Q(5,9)$.
    \vspace{0.1cm}
    \begin{tabularx}{\linewidth}{@{}c >{\raggedright\arraybackslash}p{3.8cm}
                                  >{\raggedright\arraybackslash}X
                                  >{\raggedright\arraybackslash}p{1.6cm}@{}}
        \toprule
        \textbf{Langkah} & \textbf{Operasi} & \textbf{Perhitungan} &
        \textbf{Hasil} \\
        \midrule
        1 & Periksa $P$ & $4^2=16\equiv 5$ dan $2^3+2+6=16\equiv 5$ &
        $P$ pada kurva \\
        \addlinespace
        2 & Gradien & $m \equiv (9-4)(5-2)^{-1} = 5 \cdot 3^{-1}
        = 5 \cdot 4 = 20$ & $m \equiv 9$ \\
        \addlinespace
        3 & Absis & $x_r \equiv 9^2 - 2 - 5 = 81 - 7 = 74$ &
        $x_r \equiv 8$ \\
        \addlinespace
        4 & Ordinat & $y_r \equiv 9(2-8) - 4 = -54 - 4 = -58$ &
        $y_r \equiv 8$ \\
        \bottomrule
    \end{tabularx}
    \vspace{0.15cm}
    \begin{exampleblock}{Hasil}
        \footnotesize
        $P + Q = (8,8)$. Pemeriksaan: $8^2 = 64 \equiv 9$ dan
        $8^3+8+6 = 526 \equiv 9 \pmod{11}$. Invers yang dipakai
        $3^{-1} \equiv 4$ sebab $3 \cdot 4 = 12 \equiv 1 \pmod{11}$.
    \end{exampleblock}
\end{frame}

\begin{frame}{Contoh: Penggandaan $2P$ pada $GF(11)$}
    \footnotesize
    Masih dengan kurva $y^2 \equiv x^3 + x + 6 \pmod{11}$ dan $P(2,4)$.
    \vspace{0.1cm}
    \begin{tabularx}{\linewidth}{@{}c >{\raggedright\arraybackslash}p{3.5cm}
                                  >{\raggedright\arraybackslash}X
                                  >{\raggedright\arraybackslash}p{1.5cm}@{}}
        \toprule
        \textbf{Langkah} & \textbf{Operasi} & \textbf{Perhitungan} &
        \textbf{Hasil} \\
        \midrule
        1 & Gradien & $m \equiv (3 \cdot 2^2 + 1)(2 \cdot 4)^{-1}
        = 13 \cdot 8^{-1} = 13 \cdot 7 = 91$ & $m \equiv 3$ \\
        \addlinespace
        2 & Absis & $x_r \equiv 3^2 - 2(2) = 9 - 4 = 5$ &
        $x_r = 5$ \\
        \addlinespace
        3 & Ordinat & $y_r \equiv 3(2-5) - 4 = -9 - 4 = -13$ &
        $y_r \equiv 9$ \\
        \bottomrule
    \end{tabularx}
    \vspace{0.15cm}
    \begin{exampleblock}{Hasil}
        \footnotesize
        $2P = (5,9)$. Titik ini sama dengan $Q$, sehingga
        $P+Q = 3P = (8,8)$.
    \end{exampleblock}
    \begin{alertblock}{Invers yang dipakai}
        \footnotesize
        $8^{-1} \equiv 7$ sebab $8 \cdot 7 = 56 \equiv 1 \pmod{11}$.
    \end{alertblock}
\end{frame}

\begin{frame}{Pelelaran $kP$ dan Orde Titik}
    \centering
    \footnotesize
    \begin{tabular}{@{}c l @{\hspace{7mm}} c l @{\hspace{7mm}} c l@{}}
        \toprule
        $k$ & $kP$ & $k$ & $kP$ & $k$ & $kP$ \\
        \midrule
        $1$  & $(2,4)$   & $6$  & $(7,2)$  & $11$ & $(5,2)$ \\
        \addlinespace
        $2$  & $(5,9)$   & $7$  & $(7,9)$  & $12$ & $(2,7)$ \\
        \addlinespace
        $3$  & $(8,8)$   & $8$  & $(3,6)$  & $13$ & $\mathcal{O}$ \\
        \addlinespace
        $4$  & $(10,9)$  & $9$  & $(10,2)$ &      & \\
        \addlinespace
        $5$  & $(3,5)$   & $10$ & $(8,3)$  &      & \\
        \bottomrule
    \end{tabular}
    \vspace{0.2cm}
    \begin{columns}[T,onlytextwidth]
        \begin{column}{0.48\textwidth}
            \begin{block}{Yang dibaca dari tabel}
                \footnotesize
                $3P = (8,8)$, sama dengan hasil penjumlahan $P+Q$.
            \end{block}
        \end{column}
        \begin{column}{0.48\textwidth}
            \begin{block}{Orde titik}
                \footnotesize
                $13P = \mathcal{O}$, jadi orde $P$ adalah $13$, sama dengan
                orde grup.
            \end{block}
        \end{column}
    \end{columns}
\end{frame}

% =====================================================================
\section{Penutup}
% =====================================================================

\begin{frame}{Yang Harus Diingat dan Pertemuan Berikutnya}
    \begin{block}{Butir yang harus diingat}
        \footnotesize
        \begin{itemize}
            \item Titik kurva ditambah $\mathcal{O}$ membentuk grup abelian;
                  operasinya adalah penjumlahan titik.
            \item Penjumlahan memakai gradien garis potong, penggandaan
                  memakai gradien garis singgung, semuanya modulo $p$.
            \item ECDLP --- mencari $k$ dari $Q=kP$ --- adalah dasar keamanan
                  ECC.
            \item ECDH menyamakan $abB$; ECEG mengirim
                  $\left[(kB),\ (P_M+kP_B)\right]$ dan membukanya dengan
                  $b(kB)$.
            \item ECC memberi keamanan setara dengan kunci yang jauh lebih
                  pendek daripada RSA.
        \end{itemize}
    \end{block}
    \vspace{0.1cm}
    \begin{exampleblock}{Pertemuan berikutnya}
        \textbf{Minggu 14} --- Knapsack, steganografi, dan steganalisis.
    \end{exampleblock}
\end{frame}

\end{document}
